\documentclass[11pt,a4paper]{article}
\usepackage{amsmath,amssymb}
\usepackage{physics}
\usepackage{booktabs}
\usepackage[margin=1in]{geometry}
\usepackage{hyperref}

\title{Exact solution for the momentum-diagonal interacting Chern insulator at quarter filling}
\author{Solution to Challenge 16}
\date{}

\title{Challenge 16: Final answer}
\author{}
\begin{document}
\maketitle
Thus the exact critical interaction strength is
\begin{equation}
\boxed{
U_c = 28 - 18\sqrt2 \approx 2.544155877
}
\end{equation}
in the units of the hopping amplitude used in the problem.

\section{Discussion}

The critical interaction \(U_c=28-18\sqrt2\) is the value at which the momentum-space single-occupation phase becomes thermodynamically stable against the formation of empty and doubly occupied \(k\)-blocks. This phase is an insulator with one electron per \(k\)-point, but it is not a conventional real-space Hubbard ferromagnet; the transition is driven by the competition between the kinetic energy gain of double occupation and the repulsive interaction within each momentum block.

If the interaction term in the original problem is interpreted literally as a momentum-diagonal density-density repulsion, this is the exact answer. If, instead, the notation is intended as a shorthand for a real-space on-site Hubbard interaction, the model is different and the present \(U_c\) does not apply.
\end{document}
